\documentclass[12pt]{article}
\usepackage[T1]{fontenc}
\usepackage[utf8]{inputenc}
\usepackage{lmodern}
\usepackage{textcomp}
\usepackage{enumitem}
\usepackage{geometry}
\geometry{margin=1in}
\begin{document}
\begin{center}
Answer Key - Geography - K-12 Level Assessment 16 \
\small (Answers are ordered exactly as in the question paper.)
\end{center}

\section*{Multiple Choice}
\begin{enumerate}[label=\arabic*.]
\item \textbf{Answer:} D
\\ \textit{Options:} A) Gravity concept; B) Distance decay; C) Complementarity; D) Movement bias
\\ \textit{Note:} The correct answer is "movement bias" because the establishment of new stores around a large shopping mall indicates a trend where economic activities and consumer behaviors are influenced by the centralized location, attracting more people and businesses to that area.
\item \textbf{Answer:} B
\\ \textit{Options:} A) Pioneer settlement to 1870-Germans, British, Scotch-Irish, Africans; B) 1870-1914-Asians; C) 1870-1914-Eastern and southern Europeans; D) 1965-present-Hispanics
\\ \textit{Note:} The period from 1870 to 1914 is primarily associated with the immigration of Europeans to the United States, while Asian immigration was more significantly restricted during this time due to various legislative barriers, making the association of Asians with this period incorrect.
\item \textbf{Answer:} C
\\ \textit{Options:} A) swidden.; B) subsistence farming.; C) pastorialism.; D) hunting and gathering.
\\ \textit{Note:} Pastoralism is defined as a subsistence strategy that relies on the herding and breeding of livestock for essential resources such as food, shelter, and clothing, making it the correct answer to the question.
\item \textbf{Answer:} D
\\ \textit{Options:} A) create sustainable food systems in urban areas.; B) provide jobs for women and children.; C) turn urban waste into a resource when utilized in safe manner.; D) replace agribusiness as the main focus of global food production.
\\ \textit{Note:} The correct answer is based on the understanding that while urban agriculture can supplement food production and provide local benefits, it cannot replace the scale and efficiency of agribusiness, which remains essential for meeting global food demands.
\item \textbf{Answer:} D
\\ \textit{Options:} A) innovation.; B) diffusion.; C) acculturation.; D) gravity.
\\ \textit{Note:} Gravity is a natural force that affects physical objects and is not a process through which cultural change occurs, unlike processes such as migration, trade, or communication.
\end{enumerate}

\section*{True / False}
\begin{enumerate}[label=\arabic*.]
\setcounter{enumi}{5}
\item \textbf{Answer:} False
\\ \textit{Options:} A) True; B) False
\\ \textit{Note:} The world systems theory categorizes countries into core, semi-periphery, and periphery nations, acknowledging a spectrum of economic development rather than a binary classification of rich or poor.
\item \textbf{Answer:} True
\\ \textit{Options:} A) True; B) False
\\ \textit{Note:} Acculturation is indeed the process by which immigrants assimilate aspects of the culture of their new country, including its values, language, and customs, as they adapt to their new environment.
\item \textbf{Answer:} False
\\ \textit{Options:} A) True; B) False
\\ \textit{Note:} Latitude and longitude coordinates are fixed geographic markers that define location regardless of human activities, thus they do not describe a functional region influenced by such activities.
\item \textbf{Answer:} True
\\ \textit{Options:} A) True; B) False
\\ \textit{Note:} World population is often concentrated on continental margins because these areas provide easier access to resources such as water, fertile land, and transportation routes essential for trade and commerce.
\item \textbf{Answer:} False
\\ \textit{Options:} A) True; B) False
\\ \textit{Note:} A functional region is defined by a central focus or activity, such as a city and its surrounding suburbs, rather than by cultural characteristics, which are more indicative of a formal region.
\end{enumerate}

\section*{Long Form Response}
\begin{enumerate}[label=\arabic*.]
\setcounter{enumi}{10}
\item \textbf{Answer:} Stage 1 of the demographic transition model is characterized by high birth rates and variable death rates. In this stage, societies often experience a lack of access to healthcare, sanitation, and education, leading to fluctuating death rates due to diseases and famine. As a result, the population grows slowly, as high birth rates are countered by high mortality rates. This situation can hinder societal development, as resources are primarily focused on basic survival rather than economic growth or infrastructure. Overall, Stage 1 represents a pre-industrial society where demographic changes are minimal, impacting the overall progress of the community.
\\ \textit{Note:} The answer correctly identifies that Stage 1 of the demographic transition model is marked by high birth rates and variable death rates, emphasizing how poor access to healthcare and education leads to a slow population growth that restricts societal development by diverting resources toward survival rather than economic advancement.
\item \textbf{Answer:} Barriers to the diffusion of innovation in geography include factors such as cultural resistance, economic limitations, and geographic distance. Cultural resistance occurs when communities are hesitant to adopt new ideas due to their traditional values or beliefs. Economic limitations can prevent access to new technologies, especially in less developed areas, while geographic distance can hinder communication and transport of innovations. In contrast, the Internet is not considered a barrier; rather, it serves as a powerful tool that facilitates the rapid spread of information and innovations across the globe, bridging distances and enabling diverse populations to connect and share ideas easily. Thus, the Internet enhances the diffusion of innovation rather than obstructing it.
\\ \textit{Note:} The answer correctly identifies cultural resistance, economic limitations, and geographic distance as barriers to innovation diffusion, while emphasizing that the Internet facilitates rather than hinders the spread of ideas and technologies.
\end{enumerate}

\vfill
\noindent\textit{End of Answer Key}
\end{document}